\section{Introduction}\label{sec:introduction}

Consider an unknown unit vector $s\in S^{d-1}$ and a stream of independent
Gaussian rows $X_j\sim N(0,I_d)$ with exact labels
$Y_j=\langle X_j,s\rangle$. If all raw observations can be retained, $d$
rows determine $s$ almost surely, because their matrix is invertible.
A learner with finite persistent memory must instead replace each observed
pair by an update of its state. The question is how this restriction changes
the number of observations needed for a prescribed angular accuracy.
Exact labels make the question particularly delicate: an individual real
label has no prescribed bit precision.

We study this problem through a conditional information bound. A finite
message is formed from one block of exact observations, and a separate,
independent projection of the signal is then revealed to the analyst.
We bound the information that remains in the message after this revelation.
The projection is an analysis variable throughout; it is not available to
the learner or to the rule forming the message.

\subsection*{Related memory--sample bounds}

Steinhardt and Duchi established memory and communication lower bounds for
noisy sparse regression, including a minimax squared-error lower bound
$\Omega(\sigma^2\min\{kd/(Bn),1\})$ in their model with sparsity $k$,
noise scale $\sigma$, and $n$ observations. Here $B$ bounds memory in bits
or average communication per observation
\cite[Theorem~2]{SteinhardtDuchi2015}.
For learning over finite fields, Raz proved that learning a parity requires
either quadratic memory or exponentially many samples, and extended the
branching-program method to classes of finite learning problems defined by
spectral conditions~\cite[Theorem~1]{Raz2016}\cite[Theorem~1]{Raz2017}.
These results distinguish persistent memory from unrestricted computation
on a newly observed sample.

A closely related continuous regression result is due to Sharan, Sidford,
and Valiant~\cite[Theorem~1 and Section~2]{SSV2019}. With Gaussian rows,
independent uniform noise in $[-2^{-d/5},2^{-d/5}]$, and at most $d^2/4$
bits, their bound is $\Omega(d\log r)$ observations for Euclidean accuracy
$\epsilon=d^{-r}$ in the stated range $r\le O(d/\log d)$. Their
$\Omega(d\log\log(1/\epsilon))$ consequence is stated for
$\epsilon\le d^{-\Omega(\log d)}$ within that range.
Their model already permits nonuniform finite-width
branching programs on real-valued data. The result here concerns exact
labels and $o(d^2)$ memory, and gives the single-logarithmic precision cost
$\Omega(d\log(1/\epsilon))$. These memory and noise hypotheses differ:
a lower bound with noisy labels alone does not imply a lower bound with
exact labels, since a learner receiving exact labels can add noise itself.

Dagan, Kur, and Shamir proved quadratic-space lower bounds for finding
an approximate null vector of independent Gaussian rows and for achieving
small empirical residual on a random-order consistent linear system
\cite[Theorems~3 and~10]{DKS2019}. Their tasks explain another way that
linear algebra can demand quadratic storage, but use different recovery
criteria and observation laws from the fixed-signal stream considered here.

Two companion papers give complementary treatments of exact observations.
\emph{Memory and precision in noiseless Gaussian regression}
\cite[Theorem~1.2]{OpenAIMemoryPrecision2026} obtains the same precision
scale for every fixed $A>0$ and memory $M\le Ad^2$, with a constant depending
on $A$, by propagating bounds on future success backward through blocks.
\emph{Replacing Gaussian observations in memory-constrained inference}
\cite[Proposition~7.1 and Lemma~7.2]{OpenAIReplacing2026} develops a
conditional-information comparison based on a finite two-point fiber
measure. The present paper instead develops the multi-replica block
comparisons stated below and proves its streaming consequence from them.

\subsection*{Conditional information from one block}

Let $\sigma_d$ denote uniform probability on $S^{d-1}$. In a block
experiment, the signal $S$ has law $p=f\sigma_d$, where $f$ is a bounded
nonnegative Borel density. Independently draw a standard Gaussian matrix
$A$, and observe $(A,AS)$. A finite-valued message $W$ is drawn from a
measurable kernel of these data, conditionally independently of $S$.
Finally draw a standard Gaussian matrix $B$ independently of the entire
experiment producing $(S,A,AS,W)$. Our principal quantity is
\[
 I(S;W\mid B,BS).
\]
All logarithms and entropies below are in nats.

\begin{theorem}[Gaussian replica block bound]\label{thm:gaussian-block}
For all sufficiently large $d$, set
\[
 k=2\lfloor d/16\rfloor,\qquad m=k,\qquad t=m+1,\qquad \ell=4k.
\]
Suppose $p=f\sigma_d$ is a probability with $0\le f\le L<\infty$.
In the block experiment above, let $A$ have $k$ rows and let $B$ have
$\ell$ rows. There is an absolute constant $C$ such that
\begin{equation}\label{eq:gaussian-block}
 I(S;W\mid B,BS)
 \le \frac{H(W)}{t}+Cd+C\log(2+\log L).
\end{equation}
The bound is uniform over the prior, the density bound $L$, and the message
kernel. The labels used to produce $W$ are the exact values $AS$.
\end{theorem}

The entropy cost is divided by a number $t$ proportional to $d$.
The weak dependence on $L$ matters when the theorem is applied to a
streaming state: conditioning a uniform signal on a state value $v$ of
probability $p_v>0$ gives a density at most $1/p_v$. Rare states may
therefore have large density bounds, even when their average information
cost is small.

The proof couples several possible signals with the same observed data.
Given $(A,AS)$, draw $S_1,\ldots,S_t$ independently from the posterior,
independently also of $W$ given the data. Each pair $(S_i,W)$ has the
original signal--message law. The replicas are dependent after the data
are forgotten: they satisfy $AS_1=\cdots=AS_t$ for the common hidden
matrix. Write $K$ for the relative entropy of their joint law from
$p^{\otimes t}$, and $K'$ for the corresponding relative entropy of
$(BS_1,\ldots,BS_t)$ from the product of its marginals, conditional on
$B$. These are their total correlations before and after projection.
Lemma~\ref{lem:replica-information} gives
\[
 tI(S;W\mid B,BS)\le H(W)+K-K'.
\]
Thus the proof must bound the \emph{loss} of correlation. Bounding $K$
alone would discard the geometric dependence still visible to the analyst.

The exact-label geometry supplies the comparison. For the difference
matrix $D=[S_2-S_1,\ldots,S_t-S_1]$, the replica likelihood contains
$\det(D^{\mathsf T}D)^{-k/2}$.
Theorem~\ref{thm:equal-label-measure} identifies the underlying finite
measure, including its normalization and density values on the constraint.
An observable test of the projected replicas supplies a matching
inverse-volume term.
Section~\ref{sec:gaussian} then bounds the remaining projection-density
ratio by a dyadic decomposition, proving Theorem~\ref{thm:gaussian-block}.

\subsection*{A streaming consequence}

The learner receives the fresh pairs $(X_j,Y_j)$ in order, cannot choose
the rows, and cannot revisit them. At each index it has at most $2^M$
persistent states. Its transition and stopping rules may be arbitrary
measurable randomized functions of the current state and current pair;
all information retained for later indices must be in the state.
It stops at some $\tau\in\{0,\ldots,T\}$ for a prescribed deterministic
horizon $T$. Its unit-vector output uses only the terminal state, the rule
at index $\tau$, and fresh randomness. Thus the final observation affects
the output only through that state. Rules may depend on the index, dimension,
and target accuracy. Shared randomness and initialization are jointly
independent of the signal and samples, and fresh random choices are
independent before use. The entire experiment is required to be jointly
measurable. Definition~\ref{def:streaming-model} records the full model and
the completed-measurability convention.

\begin{theorem}[Uniform-prior sample lower bound]\label{thm:streaming}
Let $M(d)=o(d^2)$ be nonnegative integer-valued, and let
$0<\epsilon(d)\le1/10$. Suppose a family of learners in
Definition~\ref{def:streaming-model} satisfies
\[
 \Prob\{\arccos\langle\widehat S,S\rangle\le\epsilon(d)\}\ge\frac23
 \qquad\text{when }S\sim\sigma_d,
\]
where the probability includes the signal, samples, and learner randomness.
Then there is an absolute $c>0$ such that
\[
 T(d)\ge c\,d\log\frac1{\epsilon(d)}
\]
for all sufficiently large $d$. The eventual threshold may depend on the
memory sequence.
\end{theorem}

The theorem assumes success averaged over the uniform sphere. A guarantee
of success at least $2/3$ for every fixed signal implies that hypothesis
by integration, so the same lower bound holds for such learners.
Section~\ref{sec:streaming-consequences} proves the result by conditioning
on the current state, applying a block estimate, and averaging its prior
cost. The independent projection leaves a residual sphere of dimension
proportional to $d$; accurate recovery still requires order
$d\log(1/\epsilon)$ conditional information on that sphere.

\subsection*{Information and geometric comparisons}

The total-correlation argument uses Watanabe's measure of multivariate
dependence~\cite{Watanabe1960} and its general probability-space chain
rules~\cite[Sections~3.3 and~6]{Austin2020}. Posterior resampling preserves
the joint law with the observed data, as in the general Nishimori identity
recorded by Lelarge and Miolane~\cite[Proposition~16]{LelargeMiolane2019}.
Here its role is to preserve the signal--message pair law while introducing
dependence among replicas. The projected tests use the
relative-entropy variational principle; a convenient statement is
\cite[Equations~(1.1)--(1.2)]{DupuisMao2022}. The equal-label measure is
related to coarea~\cite[Theorem~3.1]{Federer1959} and to affine
simplex-volume changes of variables~\cite[Lemma~1]{Drury1984}.
We prove the Gaussian probability normalization and the required
evaluations at constrained labels directly. Thus these classical tools
provide the background for the exact-fiber calculation, rather than an
unproved formula for its particular probability law.

The subsequent comparisons retain different information from that geometry.
For a prior density at most $e^b$, Theorem~\ref{thm:distance-block} uses
$q=\lfloor d/8\rfloor$ replicas and gives the alternative cost
$H(W)/q+C(d+b/d)$. Its test records successive simplex heights in distance
bins and compares the observed bin with a mixture over fixed bins.
The direct linear dependence on $b$ makes averaging over state values
immediate. The other routes use different comparison laws:
Table~\ref{tab:comparisons} records their geometric inputs and the state
information fixed during iteration. Each route verifies its own
normalization and exceptional sets. They offer distinct tests and ways to
average the prior cost, rather than a succession of stronger asymptotic
sample bounds.

\begin{table}[htbp]
\centering
\small
\setlength{\tabcolsep}{4pt}
\renewcommand{\arraystretch}{1.2}
\begin{tabular}{@{}p{.18\textwidth}p{.24\textwidth}p{.34\textwidth}p{.15\textwidth}@{}}
\toprule
Comparison & Geometric identity & Projected comparison & Conditioning \\
\midrule
Gaussian selector & Spherical equal-label measure &
Normalized determinant likelihood & Current state \\
Distance bins & Spherical equal-label measure &
Fixed mixture of height bins & Current state \\
Haar frame & Successive Haar height insertion &
Distance distribution functions & Current state \\
Synthetic Gaussian & Spherical and Gaussian equal-label measures &
Projected-prior likelihood with random covariance & Full history \\
Stiefel incidence & Regular-incidence coarea &
Conditional product law and density ratio & Full history \\
\bottomrule
\end{tabular}
\caption{The five block comparisons. The last column records the information
conditioned on before applying a block estimate. Every streaming proof also
uses a residual-sphere endpoint as in Section~\ref{sec:model-endpoint}.}
\label{tab:comparisons}
\end{table}

One complete route occupies Sections~\ref{sec:model-endpoint}--\ref{sec:streaming-consequences}.
Section~\ref{sec:model-endpoint} supplies the learner reductions and the
residual-sphere information requirement. Section~\ref{sec:replicas}
establishes the replica inequality, exact equal-label measure, and
logarithmic integrability. Sections~\ref{sec:gaussian}
and~\ref{sec:distance-bins} prove the two block estimates, and
Section~\ref{sec:streaming-consequences} gives their current-state
iterations. The final three sections develop the other geometric
comparisons and their streaming applications; the synthetic and incidence
iterations condition on full state histories. All applications use the
same exact-row, finite-state, deterministic-horizon learner model.
